\documentclass[11pt, a4paper]{article}

% --- Fonts and engine ---
\usepackage{fontspec}
\setmainfont{Helvetica Neue}
\setmonofont{Menlo}
\usepackage{unicode-math}
\setmathfont{STIX Two Math}

% --- Layout ---
\usepackage[margin=2cm, top=2cm, bottom=2.5cm]{geometry}
\usepackage{microtype}
\usepackage{parskip}
\setlength{\parskip}{0.5em}
\setlength{\parindent}{0pt}

% --- Typography and links ---
\usepackage[colorlinks=true, linkcolor=NavyBlue, urlcolor=NavyBlue, citecolor=NavyBlue]{hyperref}
\usepackage{xcolor}
\usepackage[dvipsnames]{xcolor}

% --- Math and tables ---
\usepackage{amsmath}
\usepackage{booktabs}
\usepackage{array}
\usepackage{tabularx}
\newcommand{\thead}[1]{\textbf{#1}}

% --- Graphics ---
\usepackage{graphicx}
\graphicspath{{figures/week20/}}

% --- Lists ---
\usepackage{enumitem}
\setlist[itemize]{leftmargin=1.3em, itemsep=0.15em, topsep=0.2em}
\setlist[enumerate]{leftmargin=1.5em, itemsep=0.15em, topsep=0.2em}

% --- Boxed callouts ---
\usepackage[most]{tcolorbox}
\tcbuselibrary{breakable, skins, theorems}

% Color palette
\definecolor{cardBg}{HTML}{FFFDF2}
\definecolor{cardBd}{HTML}{D4C69A}
\definecolor{defBg}{HTML}{EAF2FB}
\definecolor{defBd}{HTML}{1F4E79}
\definecolor{factBg}{HTML}{E7F1FF}
\definecolor{factBd}{HTML}{1D4ED8}
\definecolor{tipBg}{HTML}{E8F5E9}
\definecolor{tipBd}{HTML}{2E7D32}
\definecolor{trapBg}{HTML}{FDE2E2}
\definecolor{trapBd}{HTML}{D23F3F}
\definecolor{pitBg}{HTML}{FFF3CD}
\definecolor{pitBd}{HTML}{D4A017}
\definecolor{quizBg}{HTML}{FFF0E8}
\definecolor{quizBd}{HTML}{B85C2A}
\definecolor{doBg}{HTML}{F5F0E0}
\definecolor{doBd}{HTML}{6A1B9A}

\newtcolorbox{card}[1][]{enhanced, breakable, colback=cardBg, colframe=cardBd, arc=4pt, boxrule=0.5pt, left=8pt, right=8pt, top=6pt, bottom=6pt, fonttitle=\bfseries\small, coltitle=black, title={#1}}
\newtcolorbox{defbox}[1][]{enhanced, breakable, colback=defBg, colframe=defBd, arc=2pt, boxrule=0.6pt, left=8pt, right=8pt, top=6pt, bottom=6pt, fonttitle=\bfseries, title={Definition: #1}}
\newtcolorbox{factbox}[1][]{enhanced, breakable, colback=factBg, colframe=factBd, arc=3pt, boxrule=0.6pt, left=8pt, right=8pt, top=5pt, bottom=5pt, fonttitle=\bfseries, title={Fact #1}}
\newtcolorbox{trapbox}[1][]{enhanced, breakable, colback=trapBg, colframe=trapBd, arc=2pt, boxrule=0.6pt, left=8pt, right=8pt, top=5pt, bottom=5pt, fonttitle=\bfseries, title={MCQ trap #1}}
\newtcolorbox{pitfallbox}[1][]{enhanced, breakable, colback=pitBg, colframe=pitBd, arc=2pt, boxrule=0.6pt, left=8pt, right=8pt, top=5pt, bottom=5pt, fonttitle=\bfseries, title={Pitfall #1}}
\newtcolorbox{quizbox}[1][]{enhanced, breakable, colback=quizBg, colframe=quizBd, arc=2pt, boxrule=0.6pt, left=8pt, right=8pt, top=5pt, bottom=5pt, fonttitle=\bfseries, title={Lecturer's in-class quiz #1}}
\newtcolorbox{dobox}[1][]{enhanced, breakable, colback=doBg, colframe=doBd, arc=2pt, boxrule=0.6pt, left=8pt, right=8pt, top=5pt, bottom=5pt, fonttitle=\bfseries, title={Worked trace #1}}

\usepackage{titlesec}
\titleformat{\section}[block]{\Large\bfseries\color{defBd}}{\thesection.}{0.6em}{}
\titleformat{\subsection}[block]{\large\bfseries\color{defBd!80!black}}{\thesubsection}{0.5em}{}
\titlespacing*{\section}{0pt}{1.3em}{0.4em}
\titlespacing*{\subsection}{0pt}{0.9em}{0.3em}

\usepackage{fancyhdr}
\pagestyle{fancy}
\fancyhf{}
\fancyhead[L]{\small CM52078 — Week 2}
\fancyhead[R]{\small Informed and Uninformed Search}
\fancyfoot[C]{\small\thepage}
\renewcommand{\headrulewidth}{0.3pt}

\newcommand{\examflag}{\textcolor{tipBd}{\textbf{[EXAM]}}\,}
\newcommand{\trapflag}{\textcolor{trapBd}{\textbf{[TRAP]}}\,}
\newcommand{\doflag}{\textcolor{doBd}{\textbf{[DO]}}\,}

\title{\vspace{-1.5cm}{\Huge\bfseries Week 2 — Informed and Uninformed Search} \\[0.3em]{\large Search problems · BFS / DFS / IDS / UCS · Greedy / A$^{*}$ · admissibility \& consistency}}
\author{\large CM52078 Classical Artificial Intelligence \,$\cdot$\, Dr Nadejda Roubtsova}
\date{}

\begin{document}

\maketitle
\thispagestyle{empty}
\vspace{-1em}

% =====================================================================
\section{Search problems — formalisation}

\subsection{Six attributes}

A search problem is fully specified by these six items; cost is optional, the other five are mandatory.

\begin{center}
\renewcommand{\arraystretch}{1.2}
\begin{tabularx}{\textwidth}{@{}p{3.3cm}X@{}}
\toprule
\thead{Attribute} & \thead{Meaning} \\
\midrule
\textbf{State space} $S$ & Set of every configuration the agent could be in. Usually defined implicitly by actions + transition model. \\
\textbf{Initial state} $s_0$ & The starting configuration, $s_0 \in S$. \\
\textbf{Goal state / test} & A specific state, or a predicate $Goal(s)$. The two are equivalent in formalism. \\
\textbf{Actions} $Actions(s)$ & Legal moves \emph{from} state $s$. State-dependent. \\
\textbf{Transition model} $Result(s,a)$ & The state you arrive at after performing $a$ in $s$. Deterministic. \\
\textbf{Action cost} $c(s,a,s')$ \emph{(optional)} & Numeric cost. Omitted $\Rightarrow$ uniform cost $\Rightarrow$ optimal $=$ fewest transitions. \\
\bottomrule
\end{tabularx}
\end{center}

\examflag{}\textbf{Formalisation drill.} Given any puzzle (8-queens, Tower of Hanoi, river crossing, sudoku, Romania map, Rubik's, 8-puzzle), fill the six rows. This is the lecturer's standard exam-style question.

\subsection{State-space graph vs search tree}

Two \emph{different} data structures. Confusing them is a guaranteed MCQ mistake.

\begin{center}
\renewcommand{\arraystretch}{1.2}
\begin{tabularx}{\textwidth}{@{}p{4.5cm}p{4.5cm}X@{}}
\toprule
\thead{} & \thead{State-space graph} & \thead{Search tree} \\
\midrule
Represents & The whole problem (all states, all transitions) & A specific exploration strategy applied to the problem \\
Each state appears & Once & Possibly many times (nodes encode \emph{paths}, not states) \\
Built & Once; defines the problem & Incrementally during search \\
Cycles & Possible & Never (it's a tree) \\
Root & N/A & Always $s_0$ \\
\bottomrule
\end{tabularx}
\end{center}

\begin{center}

\footnotesize\textit{Romania state-space graph (left) and a partial search tree from Arad (right). Arad reappears under Sibiu in the tree because that node encodes the path Arad$\to$Sibiu$\to$Arad. R\&N AIMA 4e.}
\end{center}

\textbf{Two species of problem:}
\begin{itemize}
  \item \emph{Tree-search problem} — underlying graph has no cycles; each state reachable by exactly one path. No need for explored set.
  \item \emph{Graph-search problem} — graph has cycles; maintain an \textbf{explored} set to avoid loops. Cost: memory.
\end{itemize}

\textbf{Path-checking} is a compromise: cheaper than full explored set; only checks whether a new node is already on the current path back to the root. Catches infinite loops without paying for full history.

\subsection{Terminology}

\begin{itemize}
  \item \textbf{Expansion} — picking a frontier node and generating \emph{all} its successors by applying every legal action. The expensive operation.
  \item \textbf{Generation} — producing \emph{one} successor during an expansion. A node is generated, then sits on the frontier until itself chosen for expansion.
  \item \textbf{Frontier} (open list) — nodes generated but not yet expanded. The selection rule from the frontier \emph{is} the algorithm.
  \item \textbf{Explored} (closed list) — states already expanded; maintained in graph search.
\end{itemize}

\subsection{Goal-test timing — \trapflag\ critical}

\begin{itemize}
  \item \textbf{BFS} may test on \emph{generation} (early goal test). Safe under uniform costs --- once the goal is generated at depth $d$, no shallower one exists.
  \item \textbf{DFS} normally tests on expansion (early test also safe).
  \item \textbf{UCS} and \textbf{A$^{*}$} \emph{must} test on \emph{expansion}, never on generation. A cheaper path may still be on the frontier. Testing on generation can return a suboptimal solution.
\end{itemize}

% =====================================================================
\section{Performance metrics}

\subsection{The four metrics}

\begin{itemize}
  \item \textbf{Completeness} --- guaranteed to find a solution if one exists, or correctly declare ``no solution''.
  \item \textbf{Optimality} --- guaranteed to return the \emph{cheapest} solution (lowest path cost; with uniform costs, fewest transitions).
  \item \textbf{Time} --- measured in node expansions.
  \item \textbf{Space} --- driven by frontier size (plus explored set if maintained).
\end{itemize}

\subsection{Tree parameters $b$, $d$, $m$}

\begin{itemize}
  \item $b$ = \textbf{branching factor} (max successors per node).
  \item $d$ = depth of the \emph{shallowest} goal node.
  \item $m$ = \emph{maximum} depth of the whole tree.
  \item Always $d \le m$. Swap them and you get the wrong complexity.
\end{itemize}

\begin{center}

\footnotesize\textit{Branching factor $b$ (green), shallowest-goal depth $d$ (red), max depth $m$ (blue).}
\end{center}

\textbf{Big-O.} $\;1 + b + b^2 + \dots + b^m = O(b^m)$. Drop everything except the dominant term.

% =====================================================================
\section{Uninformed search}

\emph{Uninformed} $=$ uses only the problem (state-space graph $\pm$ action costs). \trapflag\ Action costs count as \emph{internal} information --- UCS uses costs and is \emph{still uninformed}.

\subsection{BFS — Breadth-First Search}

\begin{card}[BFS card]
\begin{itemize}
  \item \textbf{Rule:} expand the shallowest frontier node (FIFO queue).
  \item \textbf{Complete:} yes (finite branching).
  \item \textbf{Optimal:} yes \emph{only if all step costs are equal}.
  \item \textbf{Time:} $O(b^d)$. \quad \textbf{Space:} $O(b^d)$.
  \item \textbf{Early goal test} legal: check on generation, saves up to one whole layer.
\end{itemize}
\end{card}

\begin{center}

\footnotesize\textit{BFS sweeps depth-by-depth. Frontier evolves $[B,C] \to [C,D,E] \to [D,E,F,G]$. R\&N AIMA 4e.}
\end{center}

\subsection{DFS — Depth-First Search}

\begin{card}[DFS card]
\begin{itemize}
  \item \textbf{Rule:} expand the deepest frontier node (LIFO stack); backtrack when stuck.
  \item \textbf{Complete:} yes on finite trees; \emph{no} on infinite/cyclic graphs without cycle detection.
  \item \textbf{Optimal:} no --- can dive past a shallow goal and return a deeper one.
  \item \textbf{Time:} $O(b^m)$. \quad \textbf{Space:} $O(bm)$ --- \emph{linear; the only selling point}.
\end{itemize}
\end{card}

\subsection{IDS — Iterative Deepening Search}

\begin{card}[IDS card]
\begin{itemize}
  \item \textbf{Rule:} repeated DFS with depth limit $\ell = 0, 1, 2, \dots$ until goal found.
  \item Each iteration = bounded DFS $\Rightarrow$ linear space.
  \item Outer loop = depth-first-then-breadth-first $\Rightarrow$ never miss a shallow goal $\Rightarrow$ optimal.
  \item \textbf{Time:} $O(b^d)$ \quad \textbf{Space:} $O(bd)$ \quad \textbf{Complete + optimal} (equal costs).
\end{itemize}
\end{card}

\textbf{FAQ:} \emph{``IDS re-explores shallow levels --- shouldn't it be slower than BFS?''} The total expansion sum
\[ d \cdot b + (d-1) \cdot b^2 + \dots + 1 \cdot b^d \]
is dominated by the last term $b^d$. BFS's sum is $b + b^2 + \dots + b^d$, also dominated by $b^d$. Big-O captures the dominant term, so both are $O(b^d)$. IDS expands a constant factor more in absolute terms but wins on space.

\subsection{UCS — Uniform-Cost Search (Dijkstra)}

\begin{card}[UCS card]
\begin{itemize}
  \item \textbf{Rule:} expand the frontier node with the lowest accumulated path cost $g(n)$.
  \item \textbf{Complete:} yes if all action costs $\ge \epsilon > 0$.
  \item \textbf{Optimal:} yes for arbitrary non-negative costs.
  \item \textbf{Goal test on expansion}, not generation --- a cheaper path may still arrive.
  \item Equals BFS when all step costs are equal.
\end{itemize}
\end{card}

\textbf{Why ``still uninformed'':} costs are part of the problem representation, not an external estimate of distance-to-goal.

% =====================================================================
\section{Informed search}

\textbf{Informed} $=$ uses a heuristic function $h(n)$, an \emph{external} estimate of cost from $n$ to the goal. Must be optimistic (cannot overestimate); standard example is straight-line distance for pathfinding.

\subsection{Greedy Best-First}

\begin{card}[Greedy card]
\begin{itemize}
  \item \textbf{Rule:} expand the frontier node with the lowest $h(n)$. Ignores $g(n)$.
  \item Complete (with cycle detection); \emph{not} optimal --- cannot backtrack after the goal is generated, because $h(\text{goal}) = 0$ wins the frontier outright.
  \item Fast \emph{when} the heuristic is good; worst-case time $O(b^m)$.
\end{itemize}
\end{card}

\begin{center}

\footnotesize\textit{Greedy from Sibiu: picks Fagaras ($h=176$) over Rimnicu Vilcea ($h=193$), reaches Bucharest in two expansions. Path cost 310 mi --- but Sibiu$\to$RV$\to$Pitesti$\to$Bucharest costs 278 mi. Greedy returned the worse path.}
\end{center}

\subsection{A$^{*}$ Search}

\begin{card}[A$^{*}$ card]
\begin{itemize}
  \item \textbf{Rule:} expand the frontier node with the lowest $f(n) = g(n) + h(n)$.
  \item Combines UCS ($g$, backward grounding) with Greedy ($h$, forward speed).
  \item \textbf{Optimal} if $h$ is admissible.
  \item \textbf{Optimally efficient} if $h$ is consistent --- no admissible algorithm using the same $h$ expands fewer nodes.
  \item Goal test on \emph{expansion}, never on generation.
\end{itemize}
\end{card}

\begin{center}

\footnotesize\textit{A$^{*}$ Sibiu$\to$Bucharest. Step 1: RV ($f=273$) beats Fagaras ($f=275$). Step 2: from RV, Pitesti ($f=277$). Step 3: Fagaras still beats Pitesti at 275 $\Rightarrow$ expand, generates Bucharest via path 1 ($f=310$). Step 4: Pitesti expanded, generates Bucharest via path 2 ($f=278$) which replaces path 1. Optimal solution returned.}
\end{center}

% =====================================================================
\section{Heuristic properties (highest MCQ density)}

\subsection{Admissibility}

\begin{defbox}[Admissible heuristic]
$h(n) \le h^*(n)$ for every state $n$, where $h^*(n)$ is the true optimal cost from $n$ to the goal. \\
\textbf{Plain English:} never overestimates.
\end{defbox}

\begin{itemize}
  \item Straight-line distance in pathfinding --- admissible (road distance $\ge$ straight line).
  \item Manhattan distance on a 4-connected grid --- admissible.
  \item $h \equiv 0$ --- admissible (under-estimates trivially). A$^{*}$ with $h\equiv 0$ degenerates to UCS.
\end{itemize}

\subsection{Consistency (monotonicity)}

\begin{defbox}[Consistent heuristic]
For every $n$, every successor $n'$, and the action cost $c = c(n,a,n')$:
\[ h(n) \;\le\; c + h(n'). \]
This is the \emph{triangle inequality}.
\end{defbox}

\textbf{Picture:} you can't get a tighter estimate of distance-to-goal by detouring through any neighbour, after paying the actual step cost.

\subsection{The implication that gets tested constantly}

\begin{factbox}
\textbf{Consistent $\Rightarrow$ Admissible}. The converse does \textbf{not} hold. \\
Contrapositive: \emph{Not admissible $\Rightarrow$ Not consistent}.
\end{factbox}

\begin{center}
\renewcommand{\arraystretch}{1.15}
\begin{tabular}{@{}cccl@{}}
\toprule
\thead{Consistent?} & \thead{Admissible?} & \thead{Possible?} & \thead{A$^{*}$ guarantee} \\
\midrule
Yes & Yes & yes & optimal + optimally efficient \\
No  & Yes & yes & optimal (tree search) \\
Yes & No  & \emph{impossible} & --- \\
No  & No  & yes & no guarantee \\
\bottomrule
\end{tabular}
\end{center}

\textbf{Consistency triangle direction (2022 Q1):} the inequality is $h(n) \le c + h(n')$. Not $<$, not $=$, not $\ge$. Strict less-than excludes the perfect-heuristic boundary case; $\ge$ flips into overestimation.

% =====================================================================
\section{Lecturer's three in-class quizzes}

The lecturer explicitly flags these as \emph{``the hardest part of the lecture''}. Each quiz pairs a heuristic with a transition rule and asks: admissible? consistent?

\subsection{Quiz 1 --- 8-puzzle with Hamming distance}

\textbf{Heuristic:} number of tiles out of place vs goal. \textbf{Transition:} one slide $=$ cost 1.

\begin{center}

\footnotesize\textit{8-puzzle state $\{7,2,4 \mid 5,\_,6 \mid 8,3,1\}$. Eight tiles misplaced $\Rightarrow$ $h = 8$.}
\end{center}

\textbf{Admissible? Yes.} If $k$ tiles are out of place you need \emph{at least} $k$ slides to place them all (each slide moves one tile). So $h \le h^*$ always.

\textbf{Consistent? Yes.} Any single slide changes the Hamming count by $\pm 1$. With $c = 1$, the inequality $h(n) \le 1 + h(n')$ holds in every case (equality when moving toward goal, strict when away).

\subsection{Quiz 2 --- Rubik's Cube with Hamming distance}

\textbf{Heuristic:} number of incorrectly-coloured stickers. \textbf{Transition:} one face rotation $=$ cost 1.

\textbf{Admissible? \trapflag\ NO.} One rotation moves \emph{8 stickers simultaneously}. If those 8 stickers were the only wrong ones and the rotation fixes them all, $h(n) = 12$ but the true cost is 1. \emph{Massively overestimates}.

\textbf{Consistent? No.} The easy path: every consistent heuristic is admissible; this one isn't admissible, so it can't be consistent (contrapositive).

\textbf{General lesson:} a ``count misplaced units'' heuristic is admissible \emph{only if every action can fix at most one such unit}.

\subsection{Quiz 3 --- Maze with Manhattan distance}

\textbf{Heuristic:} $h(s) = |x_1 - x_2| + |y_1 - y_2|$. \textbf{Transition:} 4-connected (no diagonals), cost 1.

\textbf{Admissible? Yes.} Manhattan is the minimum number of grid moves between two cells in 4-connected movement, ignoring walls. Walls can only \emph{lengthen} the path. So $h \le h^*$.

\textbf{Consistent? Yes.} One step changes Manhattan distance by $\pm 1$ with cost 1, so $h(n) \le 1 + h(n')$ always holds.

\begin{trapbox}[: same heuristic, different verdict]
The \emph{same} heuristic (Hamming, Manhattan) is admissible in one transition system and not in another. Always check what one action can change.

\textbf{King chess heuristic (Worksheet Q6g):} 8-connected (diagonals allowed, cost 1). A diagonal step reduces Manhattan by 2 with cost 1 $\Rightarrow$ over-estimates $\Rightarrow$ \emph{not admissible}.
\end{trapbox}

% =====================================================================
\section{Worked traces ([DO] practice)}

Every trace below maps directly to a past-paper MCQ. Reproduce them on paper.

\subsection{\doflag\ Trace 1 --- Family-tree BFS counting (2023 Q1 / Worksheet Q1)}

Each person has 2 parents. Closest matching ancestor at generation 5. Tree extends to generation 7. Find the \emph{minimum} names checked.

BFS visits levels 1--4 in full: $2 + 4 + 8 + 16 = 30$. Then at level 5, the minimum is when the match is the \emph{first} name checked $\Rightarrow$ $30 + 1 = \mathbf{31}$.

\trapflag\ Maximum would be $30 + 32 = 62$. ``Minimum'' vs ``maximum'' is the entire question.

\subsection{\doflag\ Trace 2 --- Possible DFS orders (2022 Q5)}

Tree: $A \to \{B, C\}$, $B \to \{D, E\}$, $C \to \{F, G\}$. Sibling chosen at random.

DFS must exhaust a subtree before backtracking; can start with $B$ or $C$ from $A$, and either child within each subtree.

\textbf{Legal:} $A,B,D,E,C,F,G$ \quad $A,B,E,D,C,F,G$ \quad $A,C,G,F,B,E,D$ \quad (and five more permutations).

\textbf{Illegal:} $A,B,C,\dots$ (BFS-style: both children before a grandchild). $A,C,B,\dots$ (jumps subtree without exhausting).

\subsection{\doflag\ Trace 3 --- Admissibility check (2025 Q2 / Worksheet Q8)}

Edges: $A{\to}B$(3), $A{\to}C$(1), $A{\to}D$(4), $B{\to}G$(2), $C{\to}B$(2), $C{\to}D$(2), $D{\to}G$(1). Goal $G$.

\textbf{Step 1 --- compute true optimal costs $h^*$:}
\begin{itemize}
  \item $h^*(B) = 2$, $h^*(D) = 1$.
  \item $h^*(C) = \min(2{+}2, 2{+}1) = 3$ (via $D$).
  \item $h^*(A) = \min(3{+}2,\, 1{+}3,\, 4{+}1) = 4$ (via $C \to D$, \trapflag\ \emph{not} direct $A \to B$).
\end{itemize}

\textbf{Step 2 --- check each $h$ row-by-row:}

\begin{center}
\renewcommand{\arraystretch}{1.15}
\begin{tabular}{@{}lcccccc@{}}
\toprule
Node & $h^*$ & H1 & H2 & H3 & H4 & H5 \\
\midrule
A & 4 & 3\,\checkmark & 6\,\textcolor{red}{\textbf{X}} & 4\,\checkmark & 3\,\checkmark & 3\,\checkmark \\
B & 2 & 1\,\checkmark & 2\,\checkmark & 2\,\checkmark & 2\,\checkmark & 4\,\textcolor{red}{\textbf{X}} \\
C & 3 & 2\,\checkmark & 4\,\textcolor{red}{\textbf{X}} & 1\,\checkmark & 3\,\checkmark & 3\,\checkmark \\
D & 1 & 1\,\checkmark & 2\,\textcolor{red}{\textbf{X}} & 1\,\checkmark & 4\,\textcolor{red}{\textbf{X}} & 1\,\checkmark \\
\bottomrule
\end{tabular}
\end{center}

\textbf{Admissible: H1, H3.}

\subsection{\doflag\ Trace 4 --- Greedy vs A$^{*}$ on P--V graph (2024 Q2 \& Q3)}

Edges: $P{\to}Q$(2), $P{\to}S$(5), $P{\to}T$(3), $Q{\to}R$(5), $R{\to}V$(2), $S{\to}V$(5), $T{\to}U$(1), $U{\to}V$(3). $h$ $=$ actions to $V$. $h$-values: $V{=}0$, $R/S/U{=}1$, $Q/T{=}2$, $P{=}3$.

\textbf{Greedy (Q2).} From $P$: pick lowest $h$ among $\{Q(2), S(1), T(2)\}$ $\Rightarrow$ $S$. From $S$: generate $V$ $\Rightarrow$ done. \textbf{Expanded: $S$ only.}

\trapflag\ Actual cost via $S$: $5 + 5 = 10$. Via $Q \to R$: $2 + 5 + 2 = 9$ (cheaper). Greedy can't see $g$.

\textbf{A$^{*}$ (Q3).} $f$ values: from $P$ generate $Q$ ($f{=}4$), $S$ ($f{=}6$), $T$ ($f{=}5$). Expand $Q$ ($f{=}4$) $\Rightarrow$ generate $R$ ($f{=}8$). Expand $T$ ($f{=}5$) $\Rightarrow$ generate $U$ ($f{=}5$). Expand $U$ ($f{=}5$) $\Rightarrow$ generate $V$ ($f{=}7$, optimal path $P{\to}T{\to}U{\to}V$ at cost 7). \textbf{Expanded: $Q$, $T$, $U$, $R$.}

\textbf{Punchline.} Greedy expanded 1 node for a path of cost 10. A$^{*}$ expanded 4 for the optimal path of cost 7.

\subsection{\doflag\ Trace 5 --- BFS / DFS / IDS on binary state-space (Worksheet Q4)}

Start $=1$, children of $k$ are $2k$ and $2k+1$. Find $k=11$ (at depth 3, child of 5).

\begin{itemize}
  \item \textbf{BFS late goal test:} $1, 2, 3, 4, 5, 6, 7, 8, 9, 10, 11$.
  \item \textbf{BFS early goal test:} $1, 2, 3, 4, 5$ (11 generated as child of 5, halts).
  \item \textbf{Depth-limited DFS, limit 3:} $1, 2, 4, 8, 9, 5, 10, 11$.
  \item \textbf{IDS:}
    \begin{itemize}
      \item $\ell=0$: $1$.
      \item $\ell=1$: $1, 2, 3$.
      \item $\ell=2$: $1, 2, 4, 5, 3, 6, 7$.
      \item $\ell=3$: $1, 2, 4, 8, 9, 5, 10, 11$.
    \end{itemize}
\end{itemize}

\trapflag\ IDS \emph{restarts} from the root every iteration, so $1$ recurs each time. Late vs early goal test changes BFS trace by a whole layer.

\subsection{\doflag\ Trace 6 --- Graph search Greedy / A$^{*}$ / UCS (Worksheet Q5)}

Goal $H$; ties broken alphabetically.

\begin{center}
\renewcommand{\arraystretch}{1.15}
\begin{tabularx}{\textwidth}{@{}lXl@{}}
\toprule
\thead{Algo} & \thead{Removed-from-frontier order} & \thead{Path / cost} \\
\midrule
Greedy & $A, B, D, H$ & $A{\to}B{\to}D{\to}H$, cost 6 (suboptimal) \\
A$^{*}$ & $A, C, B, G, H$ & $A{\to}C{\to}G{\to}H$, cost 4 (optimal) \\
UCS & $A, C, B, G, F, H$ & $A{\to}C{\to}G{\to}H$, cost 4 (optimal) \\
\bottomrule
\end{tabularx}
\end{center}

A$^{*}$ skips $F$ that UCS expands --- the heuristic prunes one node. Both find the optimal path; UCS expands more.

% =====================================================================
\section{Past-paper traps consolidated}

\begin{itemize}
  \item \trapflag\ \textbf{BFS optimality requires equal step costs.} ``BFS is always optimal'' --- false. Use UCS for unequal costs.
  \item \trapflag\ \textbf{BFS space is $O(b^d)$, not $O(bd)$.} Linear space belongs to DFS / IDS.
  \item \trapflag\ \textbf{DFS upper bound is $O(b^m)$.} Worst case it explores the whole tree to depth $m$. $b^d$ is BFS's complexity.
  \item \trapflag\ \textbf{$d$ vs $m$ swap.} $d$ = shallowest goal. $m$ = deepest leaf.
  \item \trapflag\ \textbf{$f = g + h$, never $g - h$.}
  \item \trapflag\ \textbf{Consistency direction:} $h(n) \le c + h(n')$. Not $<$, not $\ge$, not $=$.
  \item \trapflag\ \textbf{Heuristic measures cost-to-goal, never cost-from-start.} A formula like $(x{-}1)+(y{-}1)$ in a maze starting at $(1,1)$ describes $g$, not $h$.
  \item \trapflag\ \textbf{Consistent $\Rightarrow$ admissible, one way only.} Mirror questions in 2023 Q3 and 2025 Q3.
  \item \trapflag\ \textbf{Optimally efficient $\neq$ finds optimal.} Two distinct claims.
  \item \trapflag\ \textbf{``DFS always expands at least as many nodes as A$^{*}$''} --- false. DFS can stumble on the optimal path immediately.
  \item \trapflag\ \textbf{BFS is complete with zero-cost actions} (depth-driven). UCS is \emph{not} (can loop forever on a zero-cost cycle).
  \item \trapflag\ \textbf{BFS not complete on infinite no-solution spaces} --- never finishes, never declares no-solution.
  \item \trapflag\ \textbf{Manhattan on king's-move chess is not admissible} (diagonals are cheap).
  \item \trapflag\ \textbf{``Minimum'' vs ``maximum''} in BFS counting matters (2023 Q1: 31 vs 62).
  \item \trapflag\ \textbf{Random sibling tie-break} gives multiple valid DFS orders. Don't restrict to left-first.
  \item \trapflag\ \textbf{Late vs early goal test} changes BFS trace by a whole layer.
  \item \trapflag\ \textbf{Greedy can be cheaper-to-search but suboptimal-in-cost.} Inverted intuition vs A$^{*}$.
  \item \trapflag\ \textbf{``Always preferable'' is essentially never true} (BFS vs DFS, etc).
\end{itemize}

% =====================================================================
\section{Algorithm comparison and decision card}

\subsection{Master table}

\begin{center}
\renewcommand{\arraystretch}{1.2}
\begin{tabularx}{\textwidth}{@{}p{1.4cm}p{2.6cm}p{1.6cm}p{1.6cm}p{1.6cm}X@{}}
\toprule
\thead{Algo} & \thead{Selection rule} & \thead{Time} & \thead{Space} & \thead{Optimal?} & \thead{Notes} \\
\midrule
BFS & shallowest & $O(b^d)$ & $O(b^d)$ & yes (equal costs) & early goal test legal \\
DFS & deepest & $O(b^m)$ & $O(bm)$ & no & only linear-space algo (uninformed) \\
IDS & DFS w/ inc.\ $\ell$ & $O(b^d)$ & $O(bd)$ & yes (equal costs) & best uninformed default \\
UCS & lowest $g(n)$ & varies & varies & yes (non-neg) & cost-aware uninformed; test on expansion \\
Greedy & lowest $h(n)$ & $O(b^m)$ worst & $O(b^m)$ worst & no & fast w/ good $h$; commits early \\
A$^{*}$ & lowest $g{+}h$ & varies & varies & yes if $h$ admissible & opt.\ efficient if $h$ consistent \\
\bottomrule
\end{tabularx}
\end{center}

\subsection{Equivalences}

\begin{itemize}
  \item BFS $=$ UCS when all step costs are equal.
  \item UCS $=$ A$^{*}$ with $h \equiv 0$.
  \item Greedy $=$ A$^{*}$ with $g$ ignored.
\end{itemize}

\subsection{``Which algorithm?''}

\begin{itemize}
  \item Equal step costs, memory cheap, want optimality $\Rightarrow$ BFS.
  \item Equal step costs, memory tight $\Rightarrow$ IDS.
  \item Unequal step costs, no heuristic $\Rightarrow$ UCS.
  \item Good admissible heuristic, want optimality $\Rightarrow$ A$^{*}$.
  \item Have heuristic, want fast (not necessarily optimal) $\Rightarrow$ Greedy.
  \item Don't care about optimality, memory tight $\Rightarrow$ DFS.
\end{itemize}

\subsection{Key formulas}

\begin{itemize}
  \item $1 + b + b^2 + \dots + b^m = O(b^m)$.
  \item A$^{*}$: $f(n) = g(n) + h(n)$.
  \item Admissibility: $h(n) \le h^*(n)$.
  \item Consistency: $h(n) \le c(n,a,n') + h(n')$.
  \item Implication: consistent $\Rightarrow$ admissible (and the contrapositive).
\end{itemize}

% =====================================================================
\section{Self-check questions (no answers)}

\textbf{Formalisation}
\begin{enumerate}
  \item Cast 8-queens, Tower of Hanoi, river crossing as search problems in the six-attribute form.
  \item Why are action costs ``internal'' information but a straight-line heuristic is ``external''?
  \item Distinguish ``state'' from ``node''. Can two tree nodes correspond to the same state?
  \item Give a problem whose state-space graph is a tree. Why is the explored set unnecessary there?
\end{enumerate}

\textbf{Metrics and uninformed search}
\begin{enumerate}\setcounter{enumi}{4}
  \item Why is BFS's time $O(b^d)$ and DFS's $O(b^m)$?
  \item If $d = m$, which uninformed algorithm has the smallest memory footprint?
  \item Give a problem where BFS is complete but DFS is not.
  \item Show that UCS reduces to BFS when all step costs are equal.
  \item Why does UCS test on expansion, not generation? Construct a graph where testing on generation returns a suboptimal path.
  \item Explain in big-O why IDS isn't slower than BFS despite re-exploring shallow levels.
  \item When would you prefer IDS over BFS? Over DFS?
\end{enumerate}

\textbf{Informed search}
\begin{enumerate}\setcounter{enumi}{11}
  \item Construct a 4-node example where Greedy fails to be optimal.
  \item What does A$^{*}$ degenerate to when $h \equiv 0$?
  \item Why must A$^{*}$ test on expansion (Sibiu--Bucharest example)?
  \item Construct an admissible heuristic for the 8-puzzle better than Hamming.
  \item For king's-move chess (8-connected, cost 1), construct an admissible heuristic.
  \item State the triangle inequality form of consistency and explain the name.
  \item Give an admissible heuristic that is not consistent.
\end{enumerate}

\textbf{Implications and traps}
\begin{enumerate}\setcounter{enumi}{18}
  \item If $h_1, h_2$ are admissible, is $\max(h_1, h_2)$ admissible? Consistent?
  \item Show consistent $\Rightarrow$ admissible.
  \item Why is ``A$^{*}$ with admissible $h$ is optimally efficient'' false in general?
  \item In graph search, what can go wrong with A$^{*}$ when $h$ is admissible but not consistent?
\end{enumerate}

\textbf{Trace-style}
\begin{enumerate}\setcounter{enumi}{22}
  \item Repeat the family-tree count for closest match at generation 6 in a tree of height 10. Minimum, average, maximum?
  \item Re-run Worksheet Q5 with $h$ values doubled. Still admissible? Still consistent? How does A$^{*}$'s order change?
  \item Run A$^{*}$ with Manhattan distance on an 8-puzzle for one step and verify $f$ is non-decreasing.
  \item Construct a graph where A$^{*}$ with consistent $h$ expands strictly fewer nodes than UCS.
\end{enumerate}

\end{document}
